\documentclass[
  jafontsize=11bp,
  jafontscale=0.92,
   jfm_yoko=jlreq,
  jlreq_notes,
]{jlreq}

%1. フォント設定
\usepackage{luatexja-fontspec}
	%1.1 欧文フォント
\setmainfont{RyuminPro-Light}
\setsansfont{FutoGoB101Pr6N-Bold}
	%1.2 和文フォント
\setmainjfont{RyuminPro-Light}
\setsansjfont{FutoGoB101Pr6N-Bold}

%2 スタイルファイル
\usepackage{sect}

%3 タイトル
\title{\vspace{-30mm}\Large \fbox{32} \sffamily 車に乗せてはみたものの…}
\date{\vspace{-15mm}}

\begin{document}

\maketitle

\sectionA{甲の罪責}
	\sectionB{}
		甲が制限速度を超過して自動車を走行させ、Aにぶつけてけがを負わせた行為につき、過失運転致傷罪（自動車運転死傷行為処罰法5条）の成否を検討する。
		
		同罪が成立するためには、「自動車の運転上必要な注意を怠」ること、すなわち、結果回避義務違反が認められることが必要である。制限速度を20キロメートル超過して走行する行為は、交通事故を発生させて人を死傷させる現実的危険性のある行為であって、結果回避義務違反が認められる。
		
		Aはけがを負っているが、甲の行為との間の因果関係は認められるか。見通しのよい田舎道において突然飛び出してきた子どもを避けることは困難であり、甲が制限速度内で気をつけて運転していたとしても、結果を回避することが十分可能であったとは直ちに評価できない。
		
		仮に因果関係が認められるとしても、Aが突然飛び出してくることを予見することは困難であり、Aの怪我を帰責できるだけの過失は認められない。
		
		したがって、甲に過失運転致傷罪は成立しない。
	\sectionB{}
	  甲がいったん保護したAを公園に降ろして逃げた行為につき、保護責任者遺棄罪（刑法218条）の成否を検討する。
	  
	  \sectionC{}
	  	Aは4歳の「幼年者」であり、さらに足の傷により歩行が困難になった「病者」であるため、自力で病院へ行くなどして自ら生命に対する危険に対処することができない要扶助者である。
	  
	  \sectionC{}
	  	甲に「保護する責任」（保護責任）は認められるか。218条が不保護のみならず遺棄についても保護責任を要求している趣旨は、要扶助者の生命が行為者に強く依存している場合に217条よりも思い処罰を基礎づけることに求められる。したがって、保護責任は、行為者が要扶助者の生命の安全を排他的に支配していると評価できるときに認められる。
	  	
	  	甲はAを他者からの救助が期待できない自車に乗せており、Aの生命の安全を排他的に支配する関係を作り出している。したがって、甲には保護責任が認められる。
	  	
	  \sectionC{}
	   「遺棄」とは、場所的離隔を生じさせることにより要扶助者を保護のない状態に置くことをいう。甲はAを車から降ろして誰もいない公園に放置しているが、これによりAの法益状態を悪化させているわけではないから不作為の遺棄（置き去り）にあたりうる。甲は保護責任者であり、当然に作為義務が認められるし、直ちに救急車を要請する、病院へ連れて行くなどの措置をとることも可能であった。
	   
	   したがって、甲の行為は遺棄にあたる。
	  
	  \sectionC{}
	  	遺棄罪を身体・生命に対する危険犯であると捉えると、処罰範囲が不当に拡大するおそれがある。また遺棄罪は条文上、危険の惹起を独立の要件としていないから、本罪は生命に対する抽象的危険犯と解すべきである。
	  	
	  	ここで、甲は、Aが転倒した際にアスファルトの道路に頭をぶつけたことを認識しており、生命に対する抽象的な危険の発生を認識していたといえる。そうであるにもかかわらず、事故の責任を追及されることを恐れてAを公園に放置しており、認容もあったといえる。したがって、甲には故意が認められる。
	  	
	  	以上より、甲には保護責任者遺棄罪（刑法218条）が成立する。
	
\sectionA{乙の罪責}
	単純遺棄罪（刑法217条）の共犯の成否を検討する。
　\sectionB{}
		乙は自動車を運転していたのでも、Aを車に乗せたわけでもないから、Aの生命の安全が乙に排他的に依存しているとはいえず、保護責任者には当たらない。

	甲に成立する保護責任者遺棄罪は身分犯であり、刑法は身分犯の共犯関係について65条で規律している。
	
	刑法65条について、1項が構成的身分の連帯作用を、2項が加減的身分の個別作用を定めたものとする見解がある。しかし、この見解は非身分者にも身分犯の共犯が成立する理由、構成的身分と加減的身分とでその取扱いが異なる理由を説明できないから妥当でない。そこで、同条1項は違法身分の法益侵害の間接的惹起に基づく連帯作用を、2項は責任判断の個別性に基づく責任身分の個別作用を規定したものと解するのが相当である。
	
	ここで、保護責任者という身分は、それによって遺棄罪の保護法益である生命に対する抽象的危険を発生させることができる地位ではなく、そのような危険の惹起に対する重い責任非難を基礎づけるものであるから、責任身分に当たる。
	
	したがって、乙には65条2項を適用して単純遺棄罪の共犯が成立しうることとなる。

\sectionB{}
	では、乙は甲の犯罪にどのような関与をしているか。乙はAを遺棄することに積極的に関与したわけではなく、甲からの問いかけに答えて甲の行為を阻止しなかったに過ぎない。したがって、共同正犯や教唆犯は成立せず、不作為の幇助犯（62条1項）が成立しうるにとどまる。
	
	ここで、自動車には甲のほかには乙しか乗っていなかったのであるから、甲の行為を阻止することができたのは乙しかいななったといえ、乙には甲の行為を阻止する作為義務が認められる。乙はこれを怠ることにより甲の犯行を容易にしており、乙はAが自動車にぶつかったことを認識しているから、故意も認められる。
	
	以上より、乙には単純遺棄罪の幇助犯（刑法62条1項、217条）が成立する。
	
\raggedleft{以上}
\end{document}